% TOP_PROBLEM: {"id": 2309001, "title": "Research Problems in Function Theory — Problem 9.1", "category_id": 8, "category": {"id": 8, "name": "analysis", "display_name": "Analysis", "description": "Limits, continuity, calculus, and function theory.", "slug": "analysis", "order_index": 8, "created_at": "2026-07-31T15:26:25.671Z"}, "status": "open"}
% FIRST_POSTED: null
% ENTRY_KIND: statement_only
% Imported from: batch07.tex; UnsolvedMath ID 2309001
\documentclass[11pt]{book}
\usepackage{fontspec}
\setmainfont{Latin Modern Roman}
\setsansfont{Latin Modern Sans}
\usepackage{amsmath,amssymb,mathtools}
\usepackage{unicode-math}
\setmathfont{Latin Modern Math}
\usepackage{xurl}
\usepackage[hidelinks]{hyperref}
\newcommand{\sref}[2]{\hyperlink{source:#1:#2}{[S#2]}}
\newcommand{\cataloguescope}[1]{\paragraph{Mathematical question.}#1}
\begin{document}
% UnsolvedMath ID 2309001; source catalogue code AMR-022-9001
\setcounter{section}{2309000}
\section{Research Problems in Function Theory — Problem 9.1}\label{2309001}
{\small\sffamily UnsolvedMath ID 2309001 · Analysis}
\subsection{Definitions and mathematical statement}

\paragraph{Shared notation.}$\mathbb N=\{1,2,\ldots\}$, and $\mathbb Z,\mathbb Q,\mathbb R,\mathbb C$ denote the integers, rationals, reals and complex numbers. Any other conventions are specified locally.


Let $\mathbb D=\{z\in\mathbb C:|z|<1\}$, and let $H^{\infty}(D)$ denote the algebra of bounded analytic functions on a planar domain $D$, with supremum norm. For a sequence $(z_j)\subset\mathbb D$, analytic interpolation means that every bounded complex sequence $(a_j)$ equals $(h(z_j))$ for some bounded analytic $h$ on $\mathbb D$. Harmonic interpolation means the analogous property for bounded complex harmonic functions (both real and imaginary parts are harmonic). For real data it is equivalent to using bounded real harmonic functions.
\paragraph{Statement recorded in the supplied source.}
Give a proof that analytic interpolation and harmonic interpolation are equivalent without using pre-existing conditions characterizing or guaranteeing analytic interpolating sequences. The equivalence itself is recorded as a known theorem. \sref{2309001}{2}
\subsection{Short English statement}
Prove directly that interpolation of arbitrary bounded data by bounded harmonic functions is equivalent to interpolation by bounded analytic functions.
\subsection{Sources}
\begin{description}
\item[\hypertarget{source:2309001:1}{S1}] ULAM.AI and the UnsolvedMath Contributors, \emph{UnsolvedMath: A Curated Collection of Open Mathematics Problems}, record 2309001 (AMR-022-9001). CC BY 4.0. \url{https://huggingface.co/datasets/ulamai/UnsolvedMath}.
\item[\hypertarget{source:2309001:2}{S2}] Walter K. Hayman and Edward F. Lingham, \emph{Research Problems in Function Theory} (2018), arXiv:1809.07200. Problem 9.1, complete original question, all following updates, and the relevant chapter definitions. \url{https://arxiv.org/abs/1809.07200}.
\end{description}
\label{end:2309001}
\end{document}
